% Section 19: Divergence is defined by a volume form, not by a metric} | Part: 3 | Depends on: see docs/REWRITE_PLAN.md | Status: todo | Original lines: 2824-2845
%!TEX root = ../../../main.tex
\subsection{Divergence is defined by a volume form, not by a metric}\label{sec:divmu}

Let $\mu$ be a nowhere-vanishing $n$-form on $M$ (a \emph{volume form}; in the Riemannian case $\mu=\vol_g=\sqrt{|g|}\,\dd x^1\wedge\cdots\wedge\dd x^n$). The \emph{divergence} of a vector field $X$ is defined by
\begin{equation}\label{eq:Lmu}
\Lie_X\mu=(\operatorname{div}_\mu X)\,\mu,
\end{equation}
that is, as the relative rate of change of volume under the flow of $X$. Since $\dd\mu=0$ (it is an $(n+1)$-form), Cartan's formula \eqref{eq:cartan} gives
\begin{equation}\label{eq:divmu}
\dd(\iota_X\mu)=\Lie_X\mu=(\operatorname{div}_\mu X)\,\mu .
\end{equation}
\emph{Coordinate formula, step by step.} Let $\mu=\sqrt{|g|}\,\dd x^1\wedge\cdots\wedge\dd x^n$ and $X=X^i\partial_i$.
\begin{enumerate}
\item Contract (antiderivation rule applied repeatedly, as in Section~\ref{sec:flux}): \[\iota_X\mu=\sqrt{|g|}\sum_i(-1)^{i-1}X^i\,\dd x^1\wedge\cdots\wedge\widehat{\dd x^i}\wedge\cdots\wedge\dd x^n.\]
\item Differentiate: in $\dd(\sqrt{|g|}X^i)\wedge\dd x^1\wedge\cdots\wedge\widehat{\dd x^i}\wedge\cdots$ only $\partial_i(\sqrt{|g|}X^i)\,\dd x^i$ survives (the other $\dd x^j$ are repeated), and moving $\dd x^i$ to its place gives another $(-1)^{i-1}$. The two signs cancel: $\dd(\iota_X\mu)=\sum_i\partial_i(\sqrt{|g|}X^i)\,\dd x^1\wedge\cdots\wedge\dd x^n$.
\item Divide by $\mu$:
\begin{equation}\label{eq:divcoord}
\operatorname{div}_\mu X=\frac1{\sqrt{|g|}}\,\partial_i\bigl(\sqrt{|g|}\,X^i\bigr),
\end{equation}
which reduces to $\partial_iX^i$ in Cartesian coordinates ($|g|=1$).
\end{enumerate}
So the divergence theorem is nothing but Theorem~\ref{thm:stokes} applied to the $(n-1)$-form $\omega=\iota_X\mu$ obtained from $X$ by (partial) contraction.
